% Generated draft: requires mathematical and source review.
\documentclass[11pt]{article}
\usepackage[margin=25mm]{geometry}
\usepackage{fontspec}
\setmainfont{Latin Modern Roman}
\usepackage{amsmath,amssymb,mathtools,unicode-math}
\setmathfont{Latin Modern Math}
\usepackage{enumitem,needspace,xurl,hyperref}
\setlength{\parindent}{0pt}
\setlength{\parskip}{5pt}
\setlength{\emergencystretch}{3em}
\newcommand{\N}{\mathbb{N}}
\newcommand{\Z}{\mathbb{Z}}
\newcommand{\Q}{\mathbb{Q}}
\newcommand{\R}{\mathbb{R}}
\newcommand{\C}{\mathbb{C}}
\newcommand{\F}{\mathbb{F}}
\newcommand{\A}{\mathbb{A}}
\newcommand{\PP}{\mathbb{P}}
\newcommand{\E}{\mathbb{E}}
\newcommand{\eps}{\varepsilon}
\newcommand{\dd}{\,\mathrm{d}}
\begin{document}
\section*{UnsolvedMath ID: 30000202}
\textbf{Source problem code:} catalogue entry
\subsection{Definitions and mathematical statement}
A Moufang set consists of a set \(X\) with \(|X|\geq3\) and root groups \(U_x\leq\operatorname{Sym}(X)\), one for each \(x\in X\), such that \(U_x\) fixes \(x\), acts sharply transitively on \(X\setminus\{x\}\), and the root-group family is invariant under conjugation by the little projective group \(G^\dagger=\langle U_x:x\in X\rangle\). In coordinates \(M(U,\tau)\), identify \(X=U\cup\{\infty\}\), use additive notation for the possibly nonabelian root group \(U\), and let \(\tau\) interchange \(0\) and \(\infty\). In the standard presentation, speciality is the condition
\[
(-u)\tau=-(u\tau)\qquad(u\in U\setminus\{0\}),
\]
where minus means group inverse. The presentation conventions must be matched to those in the source.

Conjecture 1: for every \(M=M(U,\tau)\) whose \(G^\dagger\)-action is not sharply two-transitive,
\[
M\text{ is special}\quad\Longleftrightarrow\quad U\text{ is abelian}.
\]
Sharply two-transitive means that there is exactly one element carrying any ordered pair of distinct points to any other such pair. The exclusion is often expressed by calling the Moufang set proper. Both directions are separate assertions.

Conjecture 2: for every special Moufang set \(M(U,\tau)\) with abelian root group, there exist a commutative field \(k\) and a quadratic Jordan division algebra \((J,Q,1)\) over \(k\) such that
\[
M(U,\tau)\cong M(J,Q,1).
\]
Here \(Q\) denotes the quadratic operator structure, avoiding confusion with the root group also denoted \(U\) in the original notation \((J,U,1)\). A quadratic Jordan division algebra is a unital quadratic Jordan algebra in which every nonzero element is invertible; its associated Moufang set has the additive root group of \(J\). No finite-dimensionality, characteristic restriction, or field-perfection hypothesis is added. Unlike Conjecture 1, the original separately stated Conjecture 2 does not explicitly exclude sharply two-transitive actions; that extra exclusion is not imported into it here.

The 2005 original passage says both directions of Conjecture 1 were unknown and no counterexamples were known to its authors. The supplied literature assessment checked 2026-08-21 instead reports the abelian-root-group implication toward speciality established for proper Moufang sets, with substantial partial results on the reverse implication and Jordan classification, but no full general resolution. These are different dated assessments. The linked finite-dimensional classification result does not by its label settle the unrestricted second conjecture.
\subsection{Short English statement}
For proper Moufang sets, ask whether speciality is equivalent to abelian root groups. Separately, ask whether every special Moufang set with abelian root group comes from a quadratic Jordan division algebra over a commutative field. The second original conjecture has no explicit properness restriction.
\subsection{Sources}
\begin{itemize}
\item Local UnsolvedMath registry, permanent ID 30000202, catalogue entry, \emph{Abelian Root Groups in Special Moufang Sets}.
\item \emph{Groups and Geometries} (2005), Oberwolfach Reports, pp. 647--721, DOI \url{https://doi.org/10.4171/owr/2005/12}, original Conjectures 1 and 2 as reproduced in the registry.
\item Supplied reference ``partial abelian-root-group theorem'': \url{https://doi.org/10.1016/j.aim.2009.10.005}.
\item Supplied reference ``recent finite-dimension classification result'': \url{https://arxiv.org/abs/2409.07445}.
\item Supplied literature triage checked 2026-08-21; no full author lists or precise partial-result hypotheses are reproduced.
\end{itemize}
\end{document}
